\section{Limitations and Monitoring}

\subsection{Evidence limitations}

\begin{itemize}
  \item The primary data are generated fixtures, not a lender's live applicants, decisions, or
    outcomes. External validity is not established.
  \item The release disposition is \texttt{characterization\_only}; customer evidence eligibility
    is false.
  \item AIR and SRG are descriptive association measures. The paper makes no causal inference.
  \item AOD and EOD are unavailable; equalized odds is not informative; ECE is not evaluated.
  \item The race surface is display-policy conditional because one group misses the predeclared
    count and share floors.
  \item LC11 covers only the configured gender and race fields, per attribute. It does not establish
    absence of bias for color, religion, national origin, marital status, age, receipt of public
    assistance, exercise of consumer-credit rights, or intersectional combinations. The paper's
    intersectional comparisons are diagnostic review surfaces, not evidence of intersectional
    absence of bias. Multi-level race encodings also require declared schema support.
  \item Quality scores describe selected diagnostic components. They do not establish predictive
    utility, privacy, full-distribution fidelity, or suitability for a decision process.
  \item The amplification advisory utility diagnostic is generally below its configured floor, and
    intrinsic utility is absent.
  \item The exact-output wall covers exact matches only. No formal differential privacy,
    membership-inference, linkage, or near-duplicate guarantee is provided.
  \item Forty robustness runs cover four configured scenarios and ten seeds. They do not exhaust
    datasets, configuration changes, dependencies, hardware, or attacks.
  \item No current native customer sizing, capacity, latency SLA, or production reliability claim
    is established by this paper.
\end{itemize}

\subsection{Deployment-specific validation}

A production evaluation should use the institution's governed data and intended configuration in a
controlled environment. It should predeclare outcomes, score provenance, protected-group handling,
reference policies, missingness treatment, minimum cell sizes, uncertainty, multiple testing,
utility and calibration endpoints, privacy tests, security controls, and escalation actions. Legal,
model-risk, information-security, privacy, and business owners should approve the plan before its
results are used for decisions.

\subsection{Monitoring triggers}

Revalidation should be triggered by a product or generator version change; a data schema,
population, target, policy, or protected-group change; drift beyond declared tolerances; dependency
or runtime-image change; evidence-schema change; failed signature/hash verification; missing
critical metric surfaces; or a change in the intended regulatory or operational use.

Recommended monitoring separates: data and population drift; branch quality; disparity and
intersectional screens; utility/calibration when enabled; privacy and exact-output controls;
evidence integrity; runtime health; and reviewer dispositions. Missing or non-informative metrics
must remain explicit states rather than being converted to passing zeros.

\subsection{Interpretive conclusion}

The v5.0.8 evidence demonstrates an auditable native dual-branch simulation pipeline, accurate
preservation of the configured gender disparity in amplification, a mechanically distinct intrinsic
policy control, deterministic provenance, and stable execution across the declared robustness plan.
It also demonstrates material limits: incomplete outcome-conditioned metrics, no calibration
surface, generally sub-floor advisory utility, and no formal privacy guarantee. Those strengths and
limits define the technical result together.
